\documentclass[10pt,a4paper]{amsart}
\usepackage[margin=19mm]{geometry}
\usepackage{amsmath,amssymb,mathtools}
\usepackage[scr=rsfs,cal=euler]{mathalfa}
\usepackage{xfrac}
\usepackage[unicode,hidelinks,bookmarks=false]{hyperref}
\newcommand{\ie}{i.e.\ }
\newcommand{\cf}{cf.\ }
\newcommand{\resp}{respectively\ }
\newcommand{\Subsection}{\subsection}
\providecommand{\nobreakdash}{\nobreak\hbox{-}}
\setlength{\emergencystretch}{2em}
\multlinegap0pt
\usepackage{bm}
\usepackage{bbm}
\usepackage{dsfont}
\usepackage{xspace}
\usepackage{cancel}
\usepackage{mathtools}
\usepackage{amsthm}
\makeatletter
\@ifundefined{theorem}{\newtheorem{theorem}{Theorem}}{}
\@ifundefined{corollary}{\newtheorem{corollary}[theorem]{Corollary}}{}
\@ifundefined{lemma}{\newtheorem{lemma}{Lemma}}{}
\@ifundefined{proposition}{\newtheorem{proposition}{Proposition}}{}
\makeatother
\title[Locally conformally K\"ahler manifolds with constant Levi-Ci]{Locally conformally K\"ahler manifolds with constant Levi-Civita or Bismut holomorphic sectional curvature}
\author{Shuwen Chen, Fangyang Zheng}
\date{Source: arXiv:2608.01893v2 (CC BY 4.0)}
\begin{document}
\maketitle

\noindent\emph{Source and licence.} Locally conformally K\"ahler manifolds with constant Levi-Civita or Bismut holomorphic sectional curvature. Version arXiv:2608.01893v2 (CC BY 4.0), distributed under the Creative Commons Attribution 4.0 International licence, as recorded in the arXiv licence metadata for this exact version and shown on the abstract page https://arxiv.org/abs/2608.01893v2. Full rights notice: licenses/arxiv-2608.01893-CC-BY-4.0.txt.\par\smallskip

\noindent\textit{Editorial scope.} Counted unit: the Proposition whose source label is \texttt{prop1} in the article (source0.tex lines 380--382, source label \texttt{prop1}), delivered together with the article's own complete original proof of it (source0.tex lines 384--404, one \texttt{proof} environment of 21 lines). No mathematics is written, repaired or re-derived: statement and proof are the article's own text line for line. Required context retained uncounted: lemma1 (lines 218--237); eq:H=c (lines 221--226); cor3 (lines 268--275); eq:Q (lines 270--274); lemma3 (lines 322--324); eq:RhR (lines 354--361); eq:Q2 (lines 372--374). Every definition, display and label of those lines is retained unchanged. Omitted from this package and not redistributed: 1--217, 227--267, 275--321, 325--353, 362--371, 375--379, 383--383, 405--758. Nothing inside a retained excerpt is omitted or altered: every retained body is delivered exactly as the article's lines are written, so no omission receipt of any kind accompanies this package. Citations: the retained excerpts carry no citation command at all, so no bibliography entry of the article is transcribed and no reference is redistributed. Reference adaptation: none was needed and none was made; every reference inside the delivered unit resolves inside it, so no unresolved reference or citation is produced. Printed slips kept literal: no printed word, symbol, hypothesis or number of the counted statement or of its proof was altered. Layout and class: the article's own class is \texttt{amsart} with packages including amsmath, amscd, amssymb, amsthm, bbm, latexsym, amsfonts, graphicx, xy, hyperref, geometry; the wrapper is the lane's pinned \texttt{amsart} editorial wrapper at 10pt on A4, which restates the theorem-like environments the retained text needs and adds the article's own macro definitions that the retained text needs (none), with extra wrapper packages (bm, bbm, dsfont, xspace, cancel, mathtools). The delivered numbering is the unit's own. Assets: no retained span contains a figure environment, a graphics-inclusion command, a PDF-page inclusion, a table or a TikZ picture; no image file is copied into this package, the package declares no asset dependency and no image file is redistributed with it. Licence: version arXiv:2608.01893v2 of this article is distributed under the Creative Commons Attribution 4.0 International licence, as recorded in the arXiv licence metadata for this exact version and shown on the abstract page https://arxiv.org/abs/2608.01893v2. A full rights notice is delivered at licenses/arxiv-2608.01893-CC-BY-4.0.txt. Characters: every retained excerpt is pure ASCII, so the delivered engine, XeLaTeX with its default Latin Modern text font, reports no missing character.
\begin{lemma}\label{lemma1}
Let \(D\) be a  metric connection on a Hermitian manifold \((M^n,g)\).  Then
we have
\begin{equation}\label{eq:H=c}
H^D\equiv c   \ \ \ \Longleftrightarrow \ \ \ \widehat R^D_{i\bar j k\bar\ell}
 =\frac{c}{2}
 \left(g_{i\bar{j}}g_{k\bar{\ell}}
      +g_{i\bar{\ell}}g_{k\bar{j}}\right), \ \ \forall \ 1\leq i,j,k,\ell \leq n,
\end{equation}
under any local frame $e=\{ e_1, \ldots , e_n\}$ of type $(1,0)$ complex vector fields on $M^n$. Here we used
\begin{equation}\label{eq:sym}
 \widehat P_{i\bar j k\bar\ell}
 =\frac14\left(
 P_{i\bar j k\bar\ell}
 +P_{k\bar j i\bar\ell}
 +P_{i\bar\ell k\bar j}
 +P_{k\bar\ell i\bar j}\right)
\end{equation}
to denote the symmetrization of a tensor $P$.
\end{lemma}

\begin{corollary} \label{cor3}
Under any local unitary frame $e$ of a Hermitian manifold $(M^n,g)$, it holds that
\begin{equation}  \label{eq:Q}
\left\{ \begin{split} \widehat{R^b} = \widehat{R} -Q,    \ \ \ \ \ \widehat{R^g} = \widehat{R} -\frac12 Q, \hspace{2.3cm}\\
  Q = \frac14 \sum_s \big( T^{j}_{is} \overline{ T^k_{\ell s }} + T^{\ell }_{ks} \overline{ T^i_{j s }} + T^{j}_{ks} \overline{ T^i_{\ell s }} + T^{\ell}_{is} \overline{ T^k_{j s }} \big).
  \end{split} \right.
\end{equation}
\end{corollary}

\begin{lemma} \label{lemma3}
A K\"ahler manifold $(M^n,g)$ is Bochner-K\"ahler if and only if there exists a Hermitian symmetric $(1,1)$ tensor $P\in {\mathcal H}$ on $M^n$ such that $R=L_g(P)$.
\end{lemma}

\begin{lemma}
Under any unitary frame $e$ for $g$ and the corresponding unitary frame $\varepsilon=e^{-u}e$ for the conformal metric $h=e^{2u}g$, their Chern torsion and curvature components are related by
\begin{eqnarray}
&& R^h_{i\bar{j}k\bar{\ell}} = e^{-2u} \big( R_{i\bar{j}k\bar{\ell}} - 2u_{i\bar{j}}  \delta_{k\ell} \big),  \label{eq:RhR}\\
&& (T^h)^j_{ik} = e^{-u} \big( T^j_{ik} + 2u_i\delta_{jk} -2 u_k \delta_{ji} \big),  \label{eq:ThT}
\end{eqnarray}
for any $1\leq i,j,k, \ell \leq n$. Here $u_i= (\partial u)(e_i)=e_i(u)$, and $u_{i\bar{j}}= (\partial \overline{\partial } u)(e_i, \overline{e}_j)$.
\end{lemma}

\begin{equation}  \label{eq:Q2}
Q =  |\nabla^g u|^2 \,G - L_g(\xi), \ \ \ \ \ \ \xi_{i\bar{j}} = u_i \overline{u_j} =(\partial u\wedge \overline{\partial}u)(e_i,\overline{e}_j).
\end{equation}

\begin{proposition}  \label{prop1}
Let $(M^n,g)$ be a Hermitian manifold, and $h=e^{2u}g$ where $u\in C^{\infty}(M)$ is a real-valued smooth function. If $h$ is K\"ahler, and $g$ has constant Chern (or Levi-Civita, or Bismut) holomorphic sectional curvature, then $h$ is Bochner-K\"ahler.
\end{proposition}

\begin{proof}
If $g$ has constant Chern holomorphic sectional curvature, $H_g=c$, then by (\ref{eq:H=c}) in Lemma \ref{lemma1}, we have $\widehat{R}=\frac{c}{2}G_g$. Since $h$ is K\"ahler, its curvature $R^h$ will obey all K\"ahler symmetries, so $R^h=\widehat{R^h}$. Now by taking the symmetrization in (\ref{eq:RhR}),  we obtain
$$ R^h_{i\bar{j}k\bar{\ell}} = e^{-2u}\frac{c}{2}(\delta_{ij}\delta_{k\ell}+\delta_{i\ell}\delta_{kj} ) -\frac12 \big( \phi_{i\bar{j}} \delta_{k\ell} + \phi_{k\bar{\ell}} \delta_{ij} + \phi_{i\bar{\ell }} \delta_{kj} + \phi_{k\bar{j}} \delta_{i\ell} \big),  $$
where $\phi_{i\bar{j}} = (\partial \overline{\partial}u)(\varepsilon_i, \overline{\varepsilon}_j)$, so $\phi$ is the Hermitian symmetric $(1,1)$ tensor corresponding to the real $(1,1)$-form $\sqrt{-1}\partial \overline{\partial}u$. That is,
\begin{equation} \label{eq:Ac}
H_g=c \ \ \ \Longrightarrow \ \ \ R^h=L_h(A), \ \ A=\frac{c}{4}e^{-2u}h - \frac12 \phi.
\end{equation}
In particular, by Lemma \ref{lemma3} we know that $h$ is Bochner-K\"ahler.

Next let us assume that $g$ has constant Bismut holomorphic sectional curvature, namely, $H^b_g=c$. Then by Lemma \ref{lemma1} we have $\widehat{R^b}=\frac{c}{2}G_g$. By (\ref{eq:Q}) in Corollary \ref{cor3} and (\ref{eq:Q2}), we obtain
$$ \widehat{R} = \widehat{R^b} + Q = \frac{c}{2}G + |\nabla^g u|^2 G - L_g(\xi), $$
where $\xi$ is the Hermitian symmetric $(1,1)$ tensor corresponding to  $\sqrt{-1}\partial u\wedge \overline{\partial}u$. By taking the symmetrization in (\ref{eq:RhR}) again  we obtain
\begin{equation}  \label{eq:Ab}
H_g^b=c \ \ \ \Longrightarrow \ \ \ R^h=L_h(A^b), \ \ A^b=\big(\frac{c}{4}e^{-2u}+\frac12 |\nabla^hu|^2\big)h - \frac12 \phi -\xi.
\end{equation}
Again by  Lemma \ref{lemma3} we know that $h$ is Bochner-K\"ahler in this case. Similarly, if the metric $g$ has constant Levi-Civita holomorphic sectional curvature $H_g^g=c$, then the same argument indicates that $h$ is Bochner-K\"ahler, and in this case we have
\begin{equation} \label{eq:Ag}
H_g^g=c \ \ \ \Longrightarrow \ \ \ R^h=L_h(A^g), \ \ A^g=\big(\frac{c}{4}e^{-2u}+\frac14 |\nabla^hu|^2\big)h - \frac12 \phi -\frac12 \xi.
\end{equation}
This completes the proof of Proposition \ref{prop1}.
\end{proof}

\end{document}
